\documentclass[11pt]{article}
\usepackage{amsmath,amssymb}
\usepackage[margin=1in]{geometry}
\title{A Gram determinant of three in a Korkine--Zolotarev basis of \(E_8\)}
\author{The Last Math Competition, conjecture 00000000354}
\date{}
\begin{document}
\maketitle

The rank-two Gram determinant in every Korkine--Zolotarev (KZ) basis of
\(E_8\) equals \(3\). This disproves the claim that the Gram determinant
chain consists of powers of two.
It also agrees with the squared Gram--Schmidt lengths tabulated by
Khandani and Esmaeili \cite[p.~154]{ke}. The argument below derives the
first two lengths directly from the coordinates of \(E_8\).

Write the lattice in the standard Euclidean coordinates as
\[
 E_8=\bigl\{x\in \mathbb Z^8\cup(\mathbb Z+\tfrac12)^8:
                         \textstyle\sum_i x_i\in 2\mathbb Z\bigr\}.
\]
The minimum nonzero squared length is \(2\). For a vector with half-integral
coordinates the eight squares are at least \(1/4\) each. A nonzero integral
vector with even coordinate sum either has at least two nonzero entries or
has one nonzero even entry. Both cases give squared length at least \(2\).
The vector \((1,-1,0,0,0,0,0,0)\) attains the bound, so any first vector
\(u\) in a KZ basis has squared length \(2\).

For any \(v\in E_8\) having nonzero projection to \(u^\perp\), choose an integer
\(k\) with \(\lvert\langle u,v\rangle-2k\rvert\le 1\). Then \(w=v-ku\) is
a nonzero lattice vector with the same orthogonal projection as \(v\). Its
projected squared length satisfies
\[
 \|w\|^2-\frac{\langle u,w\rangle^2}{\|u\|^2}
        \ge 2-\frac12=\frac32.
\]
For every choice of \(u\), there is a lattice vector \(v\) with
\(\|v\|^2=2\) and \(\langle u,v\rangle=1\). If \(u\) is integral, it has
exactly two nonzero coordinates, each equal to \(1\) or \(-1\). Pick a nonzero
coordinate \(i\) and a zero coordinate \(j\), and take \(v=u_i e_i+e_j\),
where \(e_i\) is the \(i\)-th coordinate unit vector. If \(u\) is
half-integral, all its coordinates equal \(\pm\tfrac12\). Take
\(v=2u_i e_i+2u_j e_j\) for any distinct \(i,j\). In either case \(v\) has
integer coordinates of even sum, squared norm two, and inner product one
with \(u\). Thus its projection attains the lower bound, and the second
vector of any KZ basis has projected squared length \(3/2\).

For concreteness, the pair \(u=(1,-1,0,0,0,0,0,0)\) and
\(v=(1,0,1,0,0,0,0,0)\) satisfies the first two minimizations. Its
Gram--Schmidt coefficient is \(1/2\), so the pair also satisfies the usual
size reduction. It extends to a KZ basis. Indeed, its span intersects \(E_8\) in
\(\mathbb Zu+\mathbb Zv\): every point in the intersection is integral because
its last five coordinates are zero, and its coefficients on \(u,v\) are
the negative of the second coordinate and the third coordinate,
respectively. At each subsequent stage the projection along the existing
rational span is a discrete lattice, since orthogonal projection onto the
complement of a rational subspace has a rational matrix with bounded
denominators. A shortest nonzero projected vector generates the
intersection of that projected lattice with its line. Adjoining a lift
therefore preserves the property that the chosen vectors generate all
lattice points in their span. Continuing to rank eight gives a basis with
the successive shortest-projection property. Subtracting earlier basis
vectors for size reduction leaves the spans and projected lengths intact.

The Gram determinant of any first pair is the product of the first two
squared orthogonalized lengths, namely \(2(3/2)=3\). For the explicit pair,
\[
 \det\begin{pmatrix}2&1\\1&2\end{pmatrix}=3.
\]
Since \(3\) is an odd integer greater than one, it is not a power of two.

\paragraph{Lean model.} In \texttt{Gram.lean}, vectors are doubled so that all
coordinates are integers. Membership means that the coordinates have the
same parity and sum to a multiple of four. In this model the squared norm
is four times its Euclidean value and the Gram determinant is sixteen
times its Euclidean value. The Lean proof verifies the first squared norm
minimum, the universal lower bound for the next positive Gram determinant,
and an attaining neighboring root for every shortest first vector. Its
conclusion \texttt{everyKZSecondGram} gives \(48/16=3\) for all pairs obeying
the first two KZ minimizations. The theorem \texttt{explicitKZPrefix}
independently verifies the displayed pair.

\begin{thebibliography}{1}
\bibitem{ke}
A.~K. Khandani and M. Esmaeili, \emph{Successive minimization of the state
complexity of the self-dual lattices using Korkin--Zolotarev reduced basis},
Designs, Codes and Cryptography \textbf{24} (2001), 145--158.
\end{thebibliography}
\end{document}
